xv6 quản lý file đang mở theo ba tầng: bảng fd riêng của mỗi process, bảng \texttt{struct file} chung của cả hệ thống, và tài nguyên bên dưới (inode, pipe hoặc thiết bị).

\textbf{\texttt{ofile[]} của \texttt{struct proc}} (\texttt{kernel/proc.h}) là mảng \texttt{NOFILE} = 16 con trỏ \texttt{struct file *}, riêng cho từng process. Chỉ số của mảng chính là số fd mà chương trình user nhận được: \texttt{ofile[fd]} trỏ tới \texttt{struct file} của fd đó, hoặc bằng 0 nếu fd chưa dùng. Theo quy ước, fd 0, 1, 2 là input, output và error chuẩn. Vì vậy mỗi process mở được tối đa 16 file cùng lúc.

\textbf{\texttt{file[]} của \texttt{ftable}} (\texttt{kernel/file.c}) là mảng \texttt{NFILE} = 100 \texttt{struct file} dùng chung cho cả hệ thống, được bảo vệ bởi spinlock \texttt{ftable.lock}. \texttt{filealloc()} tìm ô có \texttt{ref == 0} và đặt \texttt{ref = 1}; \texttt{filedup()} tăng \texttt{ref}; \texttt{fileclose()} giảm \texttt{ref} và chỉ thật sự đóng file khi \texttt{ref} về 0. Mỗi \texttt{struct file} (\texttt{kernel/file.h}) có các trường:
\begin{itemize}
    \item \texttt{type}: \texttt{FD\_NONE} (ô trống), \texttt{FD\_PIPE}, \texttt{FD\_INODE} (file hoặc thư mục trên đĩa) hoặc \texttt{FD\_DEVICE} (thiết bị, ví dụ console);
    \item \texttt{ref}: số con trỏ trong các \texttt{ofile[]} đang trỏ tới nó;
    \item \texttt{readable}, \texttt{writable}: quyền đọc, ghi, lấy từ chế độ mở;
    \item \texttt{pipe}, \texttt{ip}: con trỏ tới \texttt{struct pipe} hoặc \texttt{struct inode} bên dưới;
    \item \texttt{off}: vị trí đọc/ghi hiện tại, dùng cho \texttt{FD\_INODE};
    \item \texttt{major}: số hiệu thiết bị, dùng để tra bảng \texttt{devsw[]} khi là \texttt{FD\_DEVICE}.
\end{itemize}

Tầng giữa quyết định cái gì được dùng chung. Hai fd trỏ tới cùng một \texttt{struct file} (sau \texttt{dup} hoặc \texttt{fork}) thì dùng chung \texttt{off}. Hai lần \texttt{open} cùng một đường dẫn cho hai \texttt{struct file} khác nhau, mỗi cái có \texttt{off} riêng, nhưng cùng trỏ tới một inode.
